% =============================================================================
%  CHAPTER 1: Document Structure & Document Classes
% =============================================================================

\chapter{Document Structure \& Document Classes}

Every \latex{} document consists of two main parts:
\begin{enumerate}
    \item \textbf{Preamble} — everything before \verb|\begin{document}|
    \item \textbf{Body} — the content between \verb|\begin{document}| and \verb|\end{document}|
\end{enumerate}

% ------------------------------------------------------------------
\section{Document Classes}
\label{sec:docclasses}

The \verb|\documentclass| command must be the very first command. It defines the type of document.

\begin{table}[H]
\centering
\begin{tabular}{lll}
\toprule
\textbf{Class} & \textbf{Purpose} & \textbf{Sectioning} \\
\midrule
\texttt{article} & Short documents, journal papers & section, subsection \\
\texttt{report}  & Longer documents, theses        & chapter, section \\
\texttt{book}    & Full books                      & part, chapter, section \\
\texttt{letter}  & Letters                         & — \\
\texttt{beamer}  & Presentations/slides            & frame \\
\texttt{scrartcl} & KOMA-Script article            & section \\
\texttt{scrreprt} & KOMA-Script report             & chapter \\
\texttt{scrbook}  & KOMA-Script book               & part, chapter \\
\bottomrule
\end{tabular}
\caption{Common document classes in \latex{}}
\end{table}

\subsection{Document Class Options}

Class options go in square brackets: \verb|\documentclass[options]{class}|

\begin{table}[H]
\centering
\begin{tabular}{ll}
\toprule
\textbf{Option} & \textbf{Effect} \\
\midrule
\texttt{12pt}        & Sets font size (10pt, 11pt, 12pt) \\
\texttt{a4paper}     & Sets paper size (a4paper, letterpaper, legalpaper) \\
\texttt{twoside}     & Formats for two-sided printing \\
\texttt{oneside}     & Formats for one-sided printing (default) \\
\texttt{landscape}   & Landscape orientation \\
\texttt{draft}       & Draft mode — marks overfull boxes \\
\texttt{final}       & Final mode (default) \\
\texttt{twocolumn}   & Two-column layout \\
\texttt{titlepage}   & Title on a separate page \\
\texttt{openright}   & Chapters start on right-hand page \\
\texttt{leqno}       & Equation numbers on left \\
\texttt{fleqn}       & Left-aligned equations \\
\bottomrule
\end{tabular}
\caption{Common document class options}
\end{table}

\begin{lstlisting}[language={[LaTeX]TeX}, caption=Basic document structure]
\documentclass[12pt, a4paper]{article}

% === PREAMBLE ===
% Load packages, define commands, set options here
\usepackage[utf8]{inputenc}
\usepackage{graphicx}
\title{My First Document}
\author{Author Name}

% === BODY ===
\begin{document}
\maketitle
Hello, World!
\end{document}
\end{lstlisting}

% ------------------------------------------------------------------
\section{Sectioning Commands}

\latex{} provides a hierarchy of sectioning commands:

\begin{lstlisting}[language={[LaTeX]TeX}, caption=Sectioning hierarchy]
\part{Part Title}          % Level -1 (book/report only)
\chapter{Chapter Title}    % Level 0  (book/report only)
\section{Section Title}    % Level 1
\subsection{Subsection}    % Level 2
\subsubsection{Sub-sub}    % Level 3
\paragraph{Paragraph}      % Level 4
\subparagraph{Subparagraph}% Level 5
\end{lstlisting}

\subsection{Starred Versions}

Each sectioning command has a starred version (e.g., \verb|\section*{Title}|) that:
\begin{itemize}
    \item Does \emph{not} add a number
    \item Does \emph{not} appear in the table of contents
\end{itemize}

% ------------------------------------------------------------------
\section{Input and Include}

Split large documents into multiple files:

\begin{lstlisting}[language={[LaTeX]TeX}, caption=File inclusion commands]
% \input{file}     -- inserts file content inline (no page break)
% \include{file}   -- inserts file with \clearpage before and after
%                    (cannot be nested)

\input{chapters/chapter1}     % Preferred for small files
\include{chapters/chapter2}   % Preferred for chapters
\end{lstlisting}

\begin{table}[H]
\centering
\begin{tabular}{lcc}
\toprule
\textbf{Feature} & \textbf{\texttt{\textbackslash input}} & \textbf{\texttt{\textbackslash include}} \\
\midrule
Page breaks    & No  & Yes \\
Can be nested  & Yes & No  \\
Selective compilation & No & Yes (\texttt{\textbackslash includeonly}) \\
\bottomrule
\end{tabular}
\caption{Comparison of \texttt{\textbackslash input} vs \texttt{\textbackslash include}}
\end{table}

% ------------------------------------------------------------------
\section{Comments}

\begin{lstlisting}[language={[LaTeX]TeX}, caption=Comments in LaTeX]
% This is a single-line comment.
% Everything after % on the line is ignored.

% For multi-line comments, use the verbatim package:
\usepackage{verbatim}
\begin{comment}
This entire block is commented out.
Line two of the comment.
Line three.
\end{comment}
\end{lstlisting}

% ------------------------------------------------------------------
\section{Special Characters}

The following characters are \important{reserved} in \latex{} and need special handling:

\begin{table}[H]
\centering
\begin{tabular}{cll}
\toprule
\textbf{Character} & \textbf{Purpose} & \textbf{To Print It} \\
\midrule
\texttt{\#}  & Macro parameter      & \verb|\#| \\
\texttt{\$}  & Math mode            & \verb|\$| \\
\texttt{\%}  & Comment              & \verb|\%| \\
\texttt{\&}  & Table column sep.    & \verb|\&| \\
\texttt{\_}  & Subscript (math)     & \verb|\_| \\
\texttt{\{}  & Group begin          & \verb|\{| \\
\texttt{\}}  & Group end            & \verb|\}| \\
\texttt{\~{}}& Non-breaking space   & \verb|\~{}| or \verb|\textasciitilde| \\
\texttt{\^{}}& Circumflex accent    & \verb|\^{}| or \verb|\textasciicircum| \\
\texttt{\textbackslash} & Command prefix & \verb|\textbackslash| \\
\bottomrule
\end{tabular}
\caption{Special characters in \latex{}}
\end{table}

\note{Accented characters can be typed directly in UTF-8: é, ü, ñ, ç, etc.}
